\subsection{Предиктивная аналитика в серверной части}

Модуль машинного обучения на Python обеспечивает высокую точность прогнозов, однако его доступность зависит от состояния отдельного сервиса: при отказе или перезапуске analytics-ml предиктивная аналитика становится недоступна. Для обеспечения отказоустойчивости серверная часть содержит собственные алгоритмы аналитики, реализованные на Rust. Эти алгоритмы работают внутри процесса backend и не зависят от внешних сервисов, что гарантирует доступность базовой аналитики в любой момент.

\textbf{Метод Хольта--Винтерса}~\cite{HoltWinters} --- тройное экспоненциальное сглаживание для прогнозирования временных рядов надоев. Модель учитывает три компоненты: уровень (level), тренд (trend) и сезонность (seasonality). Рекуррентные формулы:

\begin{align}
l_t &= \alpha \cdot (y_t - s_{t-m}) + (1 - \alpha) \cdot (l_{t-1} + b_{t-1}) \\
b_t &= \beta \cdot (l_t - l_{t-1}) + (1 - \beta) \cdot b_{t-1} \\
s_t &= \gamma \cdot (y_t - l_t) + (1 - \gamma) \cdot s_{t-m}
\end{align}

где $l_t$ --- уровень, $b_t$ --- тренд, $s_t$ --- сезонная компонента, $m$ --- период сезонности (7 дней), $\alpha$, $\beta$, $\gamma$ --- параметры сглаживания. Прогноз на $h$ шагов вперёд:

$$\hat{y}_{t+h} = l_t + h \cdot b_t + s_{t+h-m}$$

Для обнаружения выбросов используется межквартильный размах (IQR): значение считается выбросом, если оно выходит за пределы $[Q_1 - 1.5 \cdot \text{IQR}, Q_3 + 1.5 \cdot \text{IQR}]$. Для обнаружения структурных сдвигов применяется t-критерий: если среднее последних 7 дней значимо отличается от среднего предыдущих 14 дней (p-value < 0.05), фиксируется структурный сдвиг.

\textbf{Кривая лактации Вуда}~\cite{LactationWood} --- лог-линейная модель, описывающая зависимость надоя от дня лактации:

$$y(t) = a \cdot t^b \cdot e^{-ct}$$

где $y(t)$ --- надой в день лактации $t$, $a$ --- масштабный коэффициент, определяющий начальный уровень надоя, $b$ --- параметр, контролирующий скорость нарастания до пика лактации, $c$ --- параметр, определяющий скорость спада после пика. Логарифмирование приводит к линейному виду:

$$\ln y(t) = \ln a + b \cdot \ln t - c \cdot t$$

Параметры оцениваются методом наименьших квадратов. Модель позволяет определять день пика лактации $t_{\max} = -b/c$ и сравнивать фактические надои с ожидаемыми для выявления отклонений.

\textbf{Индекс здоровья} --- составной показатель на основе z-оценок:

$$\text{Health Index} = -\frac{1}{n} \sum_{i=1}^{n} \left|z_i\right|$$

где $n = 4$ --- количество показателей, $z_i = (x_i - \mu_i) / \sigma_i$ --- стандартизованное отклонение $i$-го показателя (надои, \acrshort{scc}, активность, жвачка) от среднего $\mu_i$ и стандартного отклонения $\sigma_i$ по стаду. Индекс принимает значения от $-\infty$ до 0, где значения ниже $-2$ указывают на потенциальные проблемы со здоровьем.

\textbf{Оптимизатор запуска} --- определяет оптимальную дату сухостойного периода на основе прогноза лактации и планового отёла. Стандартный период запуска --- 60 дней до отёла. Оптимизатор корректирует дату с учётом текущего надоя.

\subsubsection{Сравнение моделей прогнозирования временных рядов}

Для повышения точности прогнозирования надоев реализована система автоматического сравнения 8 моделей прогнозирования. Исторические данные (90~дней) разбиваются на обучающую (80~процентов) и тестовую (20~процентов) выборки хронологически.

\textbf{MAPE} (Mean Absolute Percentage Error) --- основная метрика качества:

$$\text{MAPE} = \frac{100\%}{n} \sum_{i=1}^{n} \frac{|y_i - \hat{y}_i|}{|y_i|}$$

где $n$ --- количество наблюдений тестовой выборки, $y_i$ --- фактическое значение надоя, $\hat{y}_i$ --- прогнозное значение. MAPE выбран как основная метрика, поскольку она инвариантна к масштабу (надои разных коров могут отличаться в 2--3 раза).

\textbf{RMSE} (Root Mean Squared Error) дополняет MAPE, показывая типичное отклонение прогноза в абсолютных единицах (кг):

$$\text{RMSE} = \sqrt{\frac{1}{n} \sum_{i=1}^{n} (y_i - \hat{y}_i)^2}$$

В отличие от MAPE, RMSE усиливает штраф за крупные ошибки благодаря возведению в квадрат, что позволяет выявлять редкие, но значительные отклонения прогноза.

Лучшей моделью признаётся та, которая демонстрирует минимальный MAPE на тестовой выборке. Реализованы следующие модели:

\textbf{1. SES} (Simple Exponential Smoothing) --- простое экспоненциальное сглаживание:

$$\hat{y}_{t+1} = \alpha \cdot y_t + (1 - \alpha) \cdot \hat{y}_t$$

где $y_t$ --- фактическое значение надоя в момент $t$, $\hat{y}_t$ --- сглаженное значение, $\alpha \in (0, 1)$ --- параметр сглаживания, оптимизируемый перебором по сетке $\{0{,}05, 0{,}10, \ldots, 0{,}95\}$ с выбором минимального SSE (суммы квадратов ошибок).

\textbf{2. SMA} (Simple Moving Average) --- простое скользящее среднее с окном $w$.

\textbf{3. WMA} (Weighted Moving Average) --- взвешенное скользящее среднее с линейно возрастающими весами:

$$\hat{y}_{t+1} = \frac{\sum_{i=1}^{w} w_i \cdot y_{t-w+i}}{\sum_{i=1}^{w} w_i}$$

где $w$ --- размер окна (по умолчанию 7 дней), $w_i = i$ --- линейно возрастающий вес, придающий большую значимость недавним наблюдениям.

\textbf{4. Linear Regression} --- линейная регрессия по времени: $\hat{y}_t = a + b \cdot t$, где $a$ --- свободный член, $b$ --- наклон тренда, оцениваемые методом наименьших квадратов. Подходит для данных с выраженным линейным трендом, но не улавливает сезонность.

\textbf{5. Holt's Linear Method} --- двойное экспоненциальное сглаживание:

\begin{align}
l_t &= \alpha \cdot y_t + (1 - \alpha) \cdot (l_{t-1} + b_{t-1}) \\
b_t &= \beta \cdot (l_t - l_{t-1}) + (1 - \beta) \cdot b_{t-1}
\end{align}

где $l_t$ --- уровень, $b_t$ --- тренд, $\alpha$ и $\beta$ --- параметры сглаживания уровня и тренда соответственно. Прогноз: $\hat{y}_{t+h} = l_t + h \cdot b_t$.

\textbf{6. Holt--Winters} --- тройное экспоненциальное сглаживание с сезонной компонентой ($m = 7$ дней).

\textbf{7. Fourier Series} --- модель на основе рядов Фурье с линейным трендом:

$$\hat{y}_t = (a + b \cdot t) + \sum_{k=1}^{K} \left[ A_k \cos\!\left(\frac{2\pi t}{T_k}\right) + B_k \sin\!\left(\frac{2\pi t}{T_k}\right) \right]$$

где $a + b \cdot t$ --- линейный тренд, $K$ --- количество гармоник (по умолчанию 3), $T_k = 7/k$ --- период $k$-й гармоники в днях, $A_k$ и $B_k$ --- коэффициенты Фурье, оцениваемые методом наименьших квадратов. Базовый период 7 дней соответствует недельной сезонности надоев.

\textbf{8. ARIMA(0,1,1)} --- модель авторегрессии интегрированного скользящего среднего:

$$\Delta \hat{y}_t = \mu + \theta \cdot \varepsilon_{t-1}$$

где $\Delta \hat{y}_t = y_t - y_{t-1}$ --- первая разность надоя, $\mu$ --- средний дрифт (среднее суточное приращение), $\theta$ --- параметр скользящего среднего, $\varepsilon_{t-1}$ --- ошибка прогноза на предыдущем шаге. Данная спецификация эквивалентна экспоненциальному сглаживанию с дрифтом; для прогноза на $h \geq 2$ шагов вклад MA-члена равен нулю, и прогноз вырождается в прямую с наклоном $\mu$.

\subsubsection{Ансамблевый прогноз}

Для повышения робастности прогнозирования реализован ансамблевый метод, объединяющий три источника прогноза: ML-модель (XGBoost), нейросетевую модель (N-BEATS~\cite{NBEATS}) из ML-сервиса и лучшую статистическую модель (Rust). Нейросетевая модель N-BEATS использует архитектуру с базисными расширениями, способную улавливать нелинейные паттерны во временных рядах, недоступные статистическим методам.

Итоговый прогноз формируется как взвешенное среднее:

$$\hat{y}_t = w_{\text{ML}} \cdot \hat{y}_t^{\text{ML}} + w_{\text{NBEATS}} \cdot \hat{y}_t^{\text{NBEATS}} + w_{\text{Rust}} \cdot \hat{y}_t^{\text{Rust}}$$

Вес каждой модели определяется обратно пропорционально её MAPE: $w_i = (1 / (\text{MAPE}_i + 1)) / \sum_j (1 / (\text{MAPE}_j + 1))$. MAPE ML-модели и N-BEATS задаётся оценочно по результатам кросс-валидации (8 и 10~процентов соответственно), MAPE Rust-модели вычисляется динамически для каждого животного. Такое взвешивание придаёт больший вес более точным моделям, сохраняя при этом вклад каждой.

Если ML-сервис недоступен, система автоматически откатывается к прогнозу лучшей Rust-модели с весом 1,0, что гарантирует доступность прогноза при любом состоянии инфраструктуры.

\subsubsection{Объяснимость прогнозов (SHAP)}

Для обеспечения интерпретируемости прогнозов ML-сервис использует метод SHAP~\cite{SHAP}, основанный на теории кооперативных игр Шепли. Для каждого прогноза вычисляются SHAP-значения, показывающие вклад каждого признака в отклонение от базового значения модели:

$$\phi_i = \sum_{S \subseteq F \setminus \{i\}} \frac{|S|!(|F|-|S|-1)!}{|F|!} \left[ f(S \cup \{i\}) - f(S) \right]$$

где $\phi_i$ --- SHAP-значение признака $i$, $F$ --- множество всех признаков, $S$ --- подмножество признаков, не содержащее $i$, $f(S)$ --- прогноз модели на подмножестве признаков $S$, а дробь определяет вес каждой коалиции по Шепли. Для XGBoost используется оптимизированный TreeExplainer, вычисляющий точные SHAP-значения за линейное время.

В интерфейсе SHAP-значения визуализируются в виде горизонтальной столбчатой диаграммы с цветовой кодировкой: зелёный для положительного вклада (увеличивает прогноз), красный для отрицательного (уменьшает). Это позволяет фермеру понять, какие именно факторы повлияли на конкретный прогноз.
